\section{Introduction et guide de lecture}\label{int:intro}

\subsection{La question}
Former un professionnel de santé prend de trois à dix ans. Une décision prise aujourd'hui sur les places de
formation, l'âge de départ, l'accueil des diplômés étrangers ou la régulation de l'installation n'agit donc sur
l'offre de soins qu'à l'horizon d'une décennie, et parfois dans d'autres territoires que ceux visés. Les projections
publiées (DREES) sont nationales ou régionales et tendancielles. Elles ne permettent pas de composer ses propres
scénarios.

Le projet construit un outil de recherche qui répond à la question suivante : \emph{quel effectif d'une profession,
dans un territoire donné, à une date donnée, sous une politique donnée ?} La réponse prend la forme d'une
distribution et non d'un chiffre unique. Elle est comparée à la tendance, et l'écart est décomposé par flux.

\subsection{Périmètre et définitions}\label{int:perimetre}
\begin{itemize}
\item \textbf{Professions.} Sages-femmes (SF), chirurgiens-dentistes (CD), pharmaciens (PH),
  masseurs-kinésithérapeutes (MK) et infirmiers (IDE). La médecine générale est exclue par le sujet. Les médecins
  spécialistes (cardiologie, dermatologie, ophtalmologie) sont prévus en phase~2 et ne sont pas modélisés.
\item \textbf{« Libéral ».} Dans tout ce document, le mot désigne le \emph{secteur privé}, lucratif ou non. Sont
  inclus :
  \begin{itemize}
  \item les libéraux et les mixtes ;
  \item les adjoints salariés d'officine ;
  \item les salariés des établissements privés.
  \end{itemize}
  Tout employeur public est exclu. Les sources emploient « libéral » au sens d'exercice indépendant ; l'écart est
  traité au chapitre~\ref{sec:donnees-ref}.
\item \textbf{Échelles.} Commune, EPCI, bassin de vie, département, région et France. Le calcul descend à la
  commune, mais l'affichage et la validation s'arrêtent au département (chapitre~\ref{sec:validation-ref}).
\item \textbf{Horizon.} 2050 est l'horizon de référence des scénarios. L'horizon testable par backtest est
  d'environ dix ans. Au-delà, les sorties sont des \emph{scénarios conditionnels}, pas des prévisions.
\item \textbf{Données.} Uniquement des données ouvertes. Chaque fichier traité porte sa source, son millésime et le
  script qui le produit (manifeste contrôlé par \texttt{v2/data/manifest.py --check}).
\end{itemize}

\subsection{L'état en un coup d'œil}\label{int:etat}
\begin{center}\small
\begin{tabular}{@{}p{0.27\textwidth}p{0.66\textwidth}@{}}
\toprule
Composant & État au 15 novembre 2026 \\
\midrule
Jumeau (chap.~\ref{sec:jumeau-ref}) & 18 critères d'acceptation sur 18\mesure, après trois revues adversariales ;
  recalcul d'un scénario en 16\,ms\mesure au département (7 professions). \\
Méthodes (chap.~\ref{sec:methodes-ref}) & C1 (cohortes) servi par le jumeau : MAPE départementale à $h=5$ de
  7,76\,\%\mesure. Budget d'erreur : 3,11 points\mesure de bruit de comptage. Intervalles recalibrés, GBM et MLP
  rejetés. \\
Validation (chap.~\ref{sec:validation-ref}) & 12 protocoles pré-enregistrés. Le grain sous le département, les
  leviers et 2050 ne sont pas validés. Dépôt de prévisions 2027 à juger en 2027. \\
Données et paramètres & 34 fichiers traités commités sous manifeste. Hyperparamètres : 29 fiches relues sur le
  texte intégral des sources le 15 novembre (chap.~\ref{sec:parametres-ref}). \\
Interface & Carte (9 leviers de politique, 2 hypothèses de contexte), phrase en langage naturel compilée de façon
  déterministe, assistant LLM ancré sur les sorties. \\
Conformité & Pas un dispositif médical ; AIPD probablement non requise ; un point réel sur les petits effectifs
  publiés (chap.~\ref{sec:conformite-ref}). \\
Module de besoin & Programme défini, non construit (chap.~\ref{sec:ouvertes-ref}). \\
\bottomrule
\end{tabular}
\end{center}

\subsection{Ce qui peut être cité, et ce qui ne doit pas l'être}\label{int:citer}
\begin{itemize}
\item \textbf{Peut être cité.}
  \begin{itemize}
  \item Les erreurs de backtest au département, horizon 1 à 5 ans, avec leur intervalle.
  \item Les écarts entre scénarios, présentés comme des effets \emph{conditionnels au modèle}.
  \item Les comptages et séparations public/privé de l'année de base, avec leur fourchette.
  \end{itemize}
\item \textbf{Ne doit pas être cité comme prévision.}
  \begin{itemize}
  \item Un effectif en 2050, y compris dans les supports de présentation du projet. La reproduction des
    projections DREES à long terme échoue (chapitre~\ref{sec:sensibilite-ref}).
  \item Un effectif communal.
  \item L'effet d'un levier présenté comme causal et mesuré.
  \end{itemize}
\end{itemize}

\subsection{Guide de lecture selon le lecteur}
\begin{itemize}
\item \textbf{Statisticien.} Méthodes (\ref{sec:methodes-ref}), puis validation (\ref{sec:validation-ref}),
  menaces à la validité (\ref{sec:menaces-ref}) et annexe~\ref{ann:incoherences}.
\item \textbf{Chercheur en modélisation, programme Apollo.} Jumeau (\ref{sec:jumeau-ref}), sensibilité
  (\ref{sec:sensibilite-ref}), questions ouvertes et module de besoin (\ref{sec:ouvertes-ref}), reproductibilité
  (\ref{sec:repro-ref}).
\item \textbf{Professionnel de la planification.} Ce chapitre, puis données (\ref{sec:donnees-ref}), interface
  (\ref{sec:produit-ref}) et conformité (\ref{sec:conformite-ref}).
\end{itemize}

\subsection{Conventions}
Chaque nombre porte un exposant de provenance :
\begin{itemize}
\item \mesure{} : mesuré sur les données ou par un script du dépôt ;
\item \derive{} : dérivé par calcul de nombres mesurés ou sourcés ;
\item \suppose{} : supposé, c'est-à-dire posé sans mesure.
\end{itemize}
Les chemins cités sont relatifs à la racine du dépôt. Quand deux fichiers du dépôt donnent des valeurs différentes
pour une même grandeur, le texte retient la mesure la plus récente, et l'annexe~\ref{ann:incoherences} liste les
deux valeurs et la raison du choix. Le rapport décrit et compare ; il ne qualifie pas les territoires (lexique de
\texttt{v2/docs/lab/lexicon.md}).
